% from ega4-2-fr.tex
\oldpage[IV-2]{116}

pour montrer que l'homomorphisme
$\Gamma(X,\mathcal F)\to\Gamma(W,\mathcal F)$ est surjectif, donc bijectif.
Or, dire que $t=f^nt'$ signifie que l'image de $t$ dans
$\Gamma(X,\mathcal F/f^n\mathcal F)$ est nulle. Mais comme
$f^k\mathcal F/f^{k+1}\mathcal F$ est isomorphe à
$\mathcal F/f\mathcal F$, donc $Z$-pur par hypothèse, on déduit de (5.10.3),
par récurrence sur $n$, que $\mathcal F/f^n\mathcal F$ est $Z$-pur. Mais par
définition l'image de $t|W=f^ns$ dans
$\Gamma(W,\mathcal F/f^n\mathcal F)$ est nulle, d'où la conclusion.

\begin{corollary}[5.10.6]
\label{IV.2.s5.10.6-fr}
Soit $\mathcal F$ un $\mathcal O_X$-Module cohérent. Pour que $\mathcal F$
soit $Z$-clos en un point $x$, il faut et il suffit que
$\operatorname{prof}_{Z_x}(\widetilde{\mathcal F}_x)\geq2$.
\end{corollary}

Cela résulte de (5.9.6) et (5.10.5).

\begin{theorem}[5.10.7]
\label{IV.2.s5.10.7-fr}
\emph{(Hartshorne).} Soient $X$ un préschéma localement noethérien, $Y$ une
partie fermée de $X$. Supposons que pour tout $y\in Y$, on ait
$\operatorname{prof}(\mathcal O_{X,y})\geq2$; alors pour toute composante
connexe $C$ de $X$, $C-(C\cap Y)$ est connexe.
\end{theorem}

On peut se borner au cas où $X$ est connexe; il résulte alors de (5.10.5) que
l'homomorphisme canonique
$\mathcal O_X\to i_*(\mathcal O_X|X-Y)$ (où $i:X-Y\to X$ est l'injection
canonique) est bijectif. Par suite l'homomorphisme de restriction
$\Gamma(X,\mathcal O_X)\to\Gamma(X-Y,\mathcal O_X)$ est aussi bijectif. Il
suffit maintenant d'appliquer le lemme (III, 7.8.6.1).

\begin{corollary}[5.10.8]
\label{IV.2.s5.10.8-fr}
Soient $X$ un préschéma localement noethérien, $d$ un entier tel que pour tout
$x\in X$, la relation $\dim(\mathcal O_x)\geq d$ entraîne
$\operatorname{prof}(\mathcal O_x)\geq2$. Supposons $X$ connexe; alors, si
$X'$, $X''$ sont deux composantes irréductibles distinctes de $X$, il existe
une suite $(X_i)_{0\leq i\leq n}$ de composantes irréductibles de $X$ telle
que $X_0=X'$, $X_n=X''$, et que, pour $1\leq i\leq n$, on ait
\[
\operatorname{codim}(X_{i-1}\cap X_i,X)\leq d-1
\]
(on dit alors que $X$ est connexe en codimension $\leq d-1$).
\end{corollary}

Si $Y$ est une partie fermée de $X$ telle que
$\operatorname{codim}(Y,X)\geq d$, on a $\dim(\mathcal O_{X,y})\geq d$ pour
tout $y\in Y$ (5.1.3), donc $\operatorname{prof}(\mathcal O_{X,y})\geq2$ pour
tout $y\in Y$, et il résulte de (5.10.7) que $X-Y$ est connexe. D'autre part,
pour que $\operatorname{codim}(Y,X)\geq d$, il faut et il suffit que pour tout
$y\in Y$, il existe un voisinage ouvert $V$ de $y$ dans $X$ tel que
$\operatorname{codim}(Y\cap V,V)\geq d$ (0, 14.2.3). Notons enfin que si
$\mathfrak F$ désigne l'ensemble des parties fermées $Y$ de $X$ de
codimension $\geq d$, la réunion de deux ensembles de $\mathfrak F$ appartient
à $\mathfrak F$ (0, 14.2.5), et tout ensemble fermé contenu dans un ensemble
de $\mathfrak F$ appartient à $\mathfrak F$, propriétés qu'on exprime encore
en disant que $\mathfrak F$ est un \emph{antifiltre} de parties fermées de
$X$. Le corollaire résulte alors du lemme topologique suivant:

\begin{lemma}[5.10.8.1]
\label{IV.2.s5.10.8.1-fr}
Soient $X$ un espace topologique connexe localement noethérien,
$\mathfrak F$ un antifiltre de parties fermées de $X$. On suppose que si $Y$
est une partie fermée de $X$ telle que pour tout $y\in Y$, il y ait un
voisinage ouvert $V$ de $y$ dans $X$ et un $Y_y\in\mathfrak F$ tels que
$V\cap Y=V\cap Y_y$, alors $Y\in\mathfrak F$. Les conditions suivantes sont
alors équivalentes:
\begin{enumerate}
\item[a)] Pour tout $Y\in\mathfrak F$, $X-Y$ est connexe.
\item[b)] Si $X'$ et $X''$ sont deux composantes irréductibles distinctes de
$X$, il existe une suite $(X_i)_{0\leq i\leq n}$ de composantes irréductibles
de $X$ telle que $X_0=X'$, $X_n=X''$ et que, pour $1\leq i\leq n$, on ait
$X_{i-1}\cap X_i\notin\mathfrak F$.
\end{enumerate}
\end{lemma}

Supposons b) vérifiée et prouvons que $U=X-Y$ est connexe pour tout
$Y\in\mathfrak F$. Si $U'$ et $U''$ sont deux composantes irréductibles
distinctes de $U$, il existe deux composantes irréductibles $X'$, $X''$ de
$X$ telles que $X'\cap U=U'$, $X''\cap U=U''$ (0$_\mathrm I$, 2.1.6);

